\chapter{Digest 栈、DigestPool 与内容加密}
\label{ch:digest-crypto}

本章详述 ebbackup 备份内核的摘要子系统（Legacy/Standard 双栈、\texttt{DigestPool} 并行化、SHA-NI 硬件路径）以及 AES-256-GCM 内容加密与 PBKDF2 密钥派生。摘要层为 Content ID 与 dedup 提供基础；加密层在存储前变换 payload，但\textbf{不改变} Content ID 语义（不变量 I-B1，见 \cref{ch:invariants}）。

\section{DigestAlgo 与 ContentHash 路由}

\subsection{双栈设计动机（v0.2）}

ebbackup 在 v0.2 引入 \texttt{DigestAlgo} 枚举，将 SHA-256/HMAC/PBKDF2 的实现分为两套\textbf{可互操作}路径：

\begin{table}[htbp]
\centering
\caption{Digest 栈对比与默认路由}
\begin{tabular}{@{}llp{5.5cm}l@{}}
\toprule
栈 & 源码 & 典型仓库 & Feature flag \\
\midrule
\texttt{kLegacy} & \filepath{digest\_legacy.cc} & \texttt{InitRepo(legacy)}、早期测试仓库 & （无） \\
\texttt{kStandard} & \filepath{digest\_standard.cc} & \texttt{eb init} 默认、C API 新仓库 & \texttt{0x08} \\
SHA-NI 加速 & \filepath{digest\_sha\_ni.cc} & \texttt{kStandard} 路径 runtime dispatch & — \\
\bottomrule
\end{tabular}
\end{table}

\begin{designbox}[ContentHash 路由语义]
\texttt{ContentHash(algo, data, len, hash\_out)} 是引擎内唯一应使用的 chunk 摘要入口（\filepath{digest.cc}）：
\begin{enumerate}[nosep]
  \item \texttt{algo == kStandard}：若 \texttt{DigestShaNiAvailable()} 则 \texttt{Sha256ShaNi}，否则 \texttt{Sha256Standard}；
  \item \texttt{algo == kLegacy}：恒 \texttt{Sha256Legacy}；
  \item \texttt{HmacSha256} / \texttt{Pbkdf2Sha256} 按 \texttt{algo} 分别路由至 Standard 或 Legacy 实现。
\end{enumerate}
\textbf{Content ID} 恒为 \texttt{SHA256(明文 content)}；digest 栈影响 HMAC/PBKDF2 与 bench 互操作细节，不改变 content hash 对同一明文的输出（\texttt{digest\_dual\_test.cc}、\texttt{digest\_standard\_test.cc} 覆盖）。
\end{designbox}

仓库级 digest 选择由 superblock \texttt{backup\_features} 决定：

\begin{lstlisting}[caption={DigestAlgo 从 superblock 解析（\filepath{digest.h}）},label={lst:digest-from-sb}]
inline bool RepoUsesStandardDigest(const BackupSuperBlock& sb) {
  return (sb.ext.backup_features & kBackupFeatureDigestStandard) != 0;
}
inline DigestAlgo DigestAlgoFromSuperBlock(const BackupSuperBlock& sb) {
  return RepoUsesStandardDigest(sb) ? DigestAlgo::kStandard
                                    : DigestAlgo::kLegacy;
}
\end{lstlisting}

\texttt{BackupEngine::Open()} 读取 superblock 后调用 \texttt{chunk\_store\_->SetDigestAlgo(digest\_algo\_)}，保证 store 层 \texttt{Put} 重算 hash 与 pipeline 一致。C API 可通过 \texttt{EB\_BACKUP\_FLAG\_LEGACY\_DIGEST} 强制 Legacy；CLI \texttt{eb init --legacy-digest} 等效。

\subsection{Legacy 与 Standard 实现差异}

两套 SHA-256 核心均为 portable C++ 实现，常量表与轮函数结构相同；差异主要体现在：

\begin{itemize}
  \item \textbf{代码路径隔离}：Legacy 供 RAR 链、bundle 加密、历史 parity 使用；Standard 为默认栈并可挂接 SHA-NI；
  \item \textbf{HMAC/PBKDF2}：\texttt{pbkdf2\_test.cc} 验证两栈对同一 password/salt/iterations 输出一致；
  \item \textbf{编译选项}：Linux Release 对 \filepath{digest\_sha\_ni.cc} 启用 \texttt{-msha}，Legacy/Standard portable 路径不依赖 SHA 扩展。
\end{itemize}

\begin{invariantbox}[I-D1：Content ID 与 digest 栈解耦]
无论 \texttt{DigestAlgo}、Pipeline 拓扑（sequential / v4 / streaming）或 CDC 实现如何，chunk 的 Content ID 均为明文 \texttt{SHA256(content)}。加密仓库存储 AES-GCM 密文，但 \texttt{Put} 在 encode/encrypt \textbf{之前}已对明文计算 hash。
\end{invariantbox}

\section{DigestPool 并行摘要（v0.5 \wave{2}，v0.9.3 增强）}

\subsection{API 与数据结构}

\textbf{实现}：\filepath{src/common/digest\_pool.cc}、\filepath{include/ebbackup/common/digest\_pool.h}

\begin{structbox}[DigestPool 公开接口]
\begin{itemize}[nosep]
  \item \texttt{DigestPool::Shared()}：进程级单例，Pipeline / FastCDC / streaming 共享；
  \item \texttt{HashRegions(algo, base, spans, span\_count, hashes\_out)}：对同一连续 \texttt{base} 缓冲区内多个 \texttt{DigestSpan\{offset,length\}} 批量计算 32B hash，输出按 span 顺序排列；
  \item \texttt{threads()}：构造时解析的 worker 数（构造后 \texttt{SetThreads} 为 no-op，线程池规模固定）。
\end{itemize}
\end{structbox}

v0.8 起采用\textbf{持久 worker 线程池}，不再 per-batch spawn；v0.9.3 引入 \textbf{intra-job span sharding}，当 \texttt{span\_count >= threads\_} 时在调用线程上临时 fork helper 线程按 span 区间并行（避免 job queue 开销）。

\subsection{线程数解析}

\begin{lstlisting}[caption={ResolveThreadCount（\filepath{digest\_pool.cc}）},label={lst:digest-threads}]
unsigned ResolveThreadCount(unsigned requested) {
  if (requested > 0) return requested;
  const char* env = std::getenv("EBBACKUP_DIGEST_THREADS");
  if (env && env[0] != '\0') {
    const long parsed = std::strtol(env, nullptr, 10);
    if (parsed > 0) return static_cast<unsigned>(parsed);
    if (parsed == 0) return 1;
  }
  const unsigned hw = std::thread::hardware_concurrency();
  const unsigned base = hw > 0 ? hw : 8;
  return std::max<unsigned>(4, std::min(16u, base / 2));
}
\end{lstlisting}

\begin{longtable}{@{}p{3.2cm}p{10cm}@{}}
\caption{\texttt{EBBACKUP\_DIGEST\_THREADS} 环境变量语义} \\
\toprule
取值 & 行为 \\
\midrule
\endfirsthead
\multicolumn{2}{c}{\tablename\ \thetable{} — 续} \\
\toprule
取值 & 行为 \\
\midrule
\endhead
未设置 & \texttt{max(4, min(16, hw/2))}；\texttt{hw==0} 时按 8 核估算 \\
正整数 $N$ & 固定 $N$ 个持久 worker \\
\texttt{0} & 强制单线程（\texttt{digest\_pool\_test} parity 路径） \\
\bottomrule
\end{longtable}

CI bench（\filepath{bench\_check.cc}、\texttt{bench\_reuse\_gate\_test.cc}）统一设 \texttt{EBBACKUP\_DIGEST\_THREADS=4}，保证 L2/L5 门禁可重复。

\subsection{HashRegions 调度策略}

\texttt{HashRegions} 按以下顺序选择执行策略：

\begin{enumerate}
  \item \textbf{Fast path}：\texttt{span\_count == 1} 或 \texttt{threads\_ <= 1} $\rightarrow$ 调用线程串行 \texttt{ContentHash}；
  \item \textbf{Span sharding}（v0.9.3）：\texttt{span\_count >= threads\_} $\rightarrow$ 按 worker 数均分 span 索引区间，spawn 临时 \texttt{std::thread}，join 后返回（无 job queue）；
  \item \textbf{Job queue}：span 数少于线程数时，将整批 span 封装为 \texttt{DigestJob} 入队，持久 worker 在 \texttt{WorkerLoop} 中 \texttt{RunJob}，调用方 \texttt{wait} on \texttt{job->done}。
\end{enumerate}

\begin{figure}[htbp]
\centering
\begin{tikzpicture}[node distance=\flowdistance]
  \node[flowstep] (in) {HashRegions 入参\\spans, span\_count};
  \node[flowdec, below=of in] (c1) {span==1\\或 threads<=1?};
  \node[flowstep, below left=1.2cm and 0.8cm of c1] (serial) {串行 ContentHash};
  \node[flowdec, below right=1.2cm and 0.8cm of c1] (c2) {span\_count\\>= threads?};
  \node[flowstep, below=of c2] (shard) {临时线程\\按 span 分片};
  \node[flowstep, right=2.4cm of shard] (queue) {DigestJob\\入 worker 队列};

  \draw[arrow] (in) -- (c1);
  \draw[arrow] (c1) -- node[left,font=\scriptsize]{是} (serial);
  \draw[arrow] (c1) -- node[right,font=\scriptsize]{否} (c2);
  \draw[arrow] (c2) -- node[left,font=\scriptsize]{是} (shard);
  \draw[arrow] (c2) -- node[above,font=\scriptsize]{否} (queue);
\end{tikzpicture}
\caption{DigestPool::HashRegions 三路调度}
\label{fig:digest-pool-routing}
\end{figure}

Pipeline \textbf{dedup-before-encode}（v0.8+）：FastCDC / streaming 先 \texttt{HashRegions} 得到 chunk hash，再 \texttt{ExistsMany} 批量查重；命中则跳过 encode/store（见 \cref{ch:chunk-store}）。

\section{SHA-NI 硬件路径（v0.7 \wave{3}）}

\subsection{CPUID 检测与 self-test}

\texttt{digest\_sha\_ni.cc} 在 x86/x64 上实现 Intel SHA Extensions 加速的 SHA-256：

\begin{enumerate}[nosep]
  \item \texttt{ProbeShaNiCpu()}：CPUID leaf 7, subleaf 0，检测 EBX bit 29（SHA）；
  \item \texttt{SelfTestShaNi()}：对消息 \texttt{"abc"} 计算 \texttt{Sha256NiCore}，与 \texttt{Sha256Standard} 逐字节比较；
  \item 结果缓存于 \texttt{g\_sha\_ni\_cached}（\texttt{memory\_order\_acquire/release}），避免重复探测；
  \item 仅当编译宏 \texttt{EBBACKUP\_SHA\_NI\_COMPILED} 且 CPUID+self-test 均通过时 \texttt{DigestShaNiAvailable()} 返回 true。
\end{enumerate}

\begin{lstlisting}[caption={ContentHash Standard 路径 dispatch},label={lst:content-hash-shani}]
void ContentHash(DigestAlgo algo, const uint8_t* data, size_t len,
                 uint8_t hash_out[32]) {
  if (algo == DigestAlgo::kStandard) {
    if (DigestShaNiAvailable()) {
      Sha256ShaNi(data, len, hash_out);
    } else {
      Sha256Standard(data, len, hash_out);
    }
  } else {
    Sha256Legacy(data, len, hash_out);
  }
}
\end{lstlisting}

L5 bench profile 打印 \texttt{sha\_ni: available|unavailable}（\filepath{backup\_pipeline.cc} phase stats）。非 x86 或未编译 SHA 内联时，\texttt{Sha256ShaNi} 透明 fallback 至 \texttt{Sha256Standard}。

\subsection{Sha256ShaNi 实现要点}

\texttt{Sha256NiCore} 使用 \texttt{\_mm\_sha256rnds2\_epu32}、\texttt{\_mm\_sha256msg1/2\_epu32} 处理 64B 块；尾部 padding 与 bit-length 编码与 FIPS 180-4 一致。\texttt{digest\_sha\_ni\_test.cc} 的 \texttt{ContentHashDispatchParity} 验证 dispatch 输出与 Standard 一致。

\section{AES-256-GCM 与 PBKDF2}

\subsection{加密管线}

\textbf{实现}：\filepath{src/crypto/aes\_gcm.cc}、\filepath{aes\_block.cc}、\filepath{gcm\_portable.cc}

\begin{table}[htbp]
\centering
\caption{AES-256-GCM 存储格式}
\begin{tabular}{@{}lll@{}}
\toprule
字段 & 长度 & 说明 \\
\midrule
nonce & 12 B & 每 chunk 随机（\texttt{kAesGcmNonceSize}） \\
ciphertext & $N$ B & 与明文等长 \\
auth tag & 16 B & GCM MAC（\texttt{kAesGcmTagSize}） \\
\bottomrule
\end{tabular}
\end{table}

Windows 优先 \texttt{BCrypt} AES-GCM；失败或非 Windows 走 \texttt{Aes256GcmEncryptPortable}。NIST 测试向量覆盖 portable 路径（\filepath{tests/crypto/}）。

\texttt{ChunkStore::PutKnownHash} 在 \texttt{ContentClassEncode} 之后、\texttt{AppendRecord} 之前调用 \texttt{Aes256GcmEncrypt}；codec 经 \texttt{EncryptCodec(codec)} 标记加密位，但 header 中 \texttt{hash[32]} 仍为\textbf{明文} Content ID。

\subsection{PBKDF2 与 crypto.salt}

\begin{lstlisting}[caption={DeriveContentKey（\filepath{aes\_gcm.cc}）},label={lst:derive-cek}]
Status DeriveContentKey(const std::string& password, const uint8_t salt[16],
                        uint8_t cek[32], DigestAlgo algo) {
  if (password.empty()) return Status::InvalidArgument("empty password");
  Pbkdf2Sha256(algo, ..., salt, 16, 100000, cek);
  return Status::Ok();
}
\end{lstlisting}

\begin{itemize}
  \item \filepath{crypto.salt}：每仓库 16B 随机 salt；\texttt{LoadOrCreateRepoSalt} 首次创建后 \texttt{FsyncPath}；
  \item PBKDF2-HMAC-SHA256，\textbf{100\,000} 次迭代，输出 32B Content Encryption Key（CEK）；
  \item digest 栈决定 PBKDF2 内部 HMAC 实现，Legacy/Standard 仓库均可加密（\texttt{digest\_dual\_encrypt\_test}）；
  \item Feature \texttt{kBackupFeatureEncrypted}（0x04）；CLI \texttt{--encrypt} / \texttt{--password-env}。
\end{itemize}

\section{ChunkRecordHeader v1 与 v2}

chunk 记录在 legacy 单文件与 EbPack 内共用 header 布局（v2 为当前默认）：

\begin{longtable}{@{}rrl p{6cm}@{}}
\caption{ChunkRecordHeader v2 字节布局（48B，\filepath{chunk\_store.h}）} \\
\toprule
偏移 & 长度 & 字段 & 说明 \\
\midrule
\endfirsthead
\multicolumn{4}{c}{\tablename\ \thetable{} — 续} \\
\toprule
偏移 & 长度 & 字段 & 说明 \\
\midrule
\endhead
0 & 32 & \texttt{hash[32]} & Content ID（明文 SHA-256） \\
32 & 4 & \texttt{stored\_len} & 存储 payload 字节数（含加密/压缩后） \\
36 & 4 & \texttt{uncompressed\_len} & 逻辑明文长度 \\
40 & 1 & \texttt{codec} & \texttt{ChunkCodec} 枚举 \\
41 & 3 & \texttt{reserved[3]} & 对齐保留 \\
44 & 4 & \texttt{record\_crc32} & header CRC（字段置零后计算） \\
\midrule
\multicolumn{4}{l}{\textbf{v1（40B）}：仅 \texttt{hash} + \texttt{raw\_len} + \texttt{record\_crc32}；无 codec/uncompressed\_len，逻辑长度等于 \texttt{raw\_len}} \\
\bottomrule
\end{longtable}

\begin{lstlisting}[caption={ChunkRecordHeader v2 结构},label={lst:chunk-header-v2}]
struct ChunkRecordHeader {
  uint8_t hash[32];
  uint32_t stored_len;
  uint32_t uncompressed_len;
  uint8_t codec;
  uint8_t reserved[3];
  uint32_t record_crc32;
};
static_assert(sizeof(ChunkRecordHeader) == 48);
\end{lstlisting}

\texttt{ReadParsedHeaderAt} 先尝试 v1 CRC；失败则按 v2 解析。Open/LoadIndex 扫描与 persistent index 均支持 v1/v2 混存仓库（向后兼容）。

\section{与 Pipeline / ChunkStore 的集成}

\begin{figure}[htbp]
\centering
\begin{tikzpicture}[node distance=\flowdistance]
  \node[flowdata] (plain) {明文 chunk};
  \node[flowstep, below=of plain] (hash) {DigestPool::HashRegions\\或 ContentHash};
  \node[flowstep, below=of hash] (dedup) {ExistsMany / Exists};
  \node[flowdec, below=of dedup] (hit) {已存在?};
  \node[flowstep, below left=1.1cm and 0.6cm of hit] (skip) {跳过 store};
  \node[flowstep, below right=1.1cm and 0.6cm of hit] (enc) {Encode + 可选 AES-GCM};
  \node[flowstep, below=of enc] (put) {PutKnownHash / PutPrecompressed};

  \draw[arrow] (plain)--(hash)--(dedup)--(hit);
  \draw[arrow] (hit)-- node[left,font=\scriptsize]{是} (skip);
  \draw[arrow] (hit)-- node[right,font=\scriptsize]{否} (enc)--(put);
\end{tikzpicture}
\caption{摘要层在 dedup-before-encode 中的位置}
\label{fig:digest-pipeline}
\end{figure}

Streaming CDC（v0.9.3）对虚拟段使用 file-absolute \texttt{digest\_base} 指针调用 \texttt{HashRegions}，保证与 \texttt{FastCdcSlice} mmap 路径 hash 一致。

\section{Legacy 栈的保留场景}

尽管 \texttt{eb init} 默认 Standard，Legacy 栈在以下场景仍为\textbf{权威}路径：

\begin{table}[htbp]
\centering
\caption{必须使用 Legacy digest 的组件}
\begin{tabular}{@{}llp{7cm}@{}}
\toprule
组件 & API & 原因 \\
\midrule
RAR 审计链 & \texttt{ComputeRarSha256(body, kLegacy)} & 链式 hash 与历史 audit 记录绑定 \\
eb bundle 加密 & \texttt{Pbkdf2Sha256(kLegacy, ...)} & 归档 bundle 互操作 \\
\texttt{Sha256()} 便捷函数 & \texttt{DigestAlgo::kLegacy} & 非引擎工具代码默认 \\
C API legacy flag & \texttt{EB\_BACKUP\_FLAG\_LEGACY\_DIGEST} & 显式兼容模式 \\
\bottomrule
\end{tabular}
\end{table}

迁移 Standard 仓库时\textbf{无需}重算 chunk hash；仅需在新 backup 时启用 \texttt{kBackupFeatureDigestStandard}，后续 HMAC/PBKDF2 走 Standard 实现。

\section{DigestPool 实现细节与边界条件}

\subsection{Worker 生命周期}

\texttt{DigestPoolImpl} 构造时 spawn \texttt{threads\_} 个 \texttt{WorkerLoop} 线程；析构时设 \texttt{stop\_=true}、\texttt{notify\_all}、\texttt{join} 全部 worker。进程退出时 \texttt{Shared()} 单例随静态析构顺序销毁——gtest 中局部 \texttt{DigestPool pool(4)} 实例会在测试结束时正常 join。

\subsection{Job queue 路径的适用场景}

当 FastCDC 一切文件仅产生少量 span（例如小文件 2--3 个 chunk）且 \texttt{threads\_>1} 时，走 job queue 而非 span sharding：单个 worker 处理整批 span，其余 worker 空闲。这是预期行为——micro-batch 的调度开销低于错误并行化带来的同步成本。大文件 streaming 路径通常 \texttt{span\_count >> threads}，命中 sharding 快路径。

\subsection{SetThreads 为 no-op 的原因}

v0.8 后线程池规模在构造时固定；\texttt{SetThreads} 保留 API 兼容性但不改变 \texttt{threads\_}。动态调整 worker 数需重建 \texttt{DigestPool} 实例（bench 通过环境变量在进程启动前设定）。

\section{SHA-NI 编译与平台矩阵}

\begin{longtable}{@{}llp{7cm}@{}}
\caption{SHA-NI 可用性矩阵} \\
\toprule
平台 & 编译 & Runtime \\
\midrule
\endfirsthead
\multicolumn{3}{c}{\tablename\ \thetable{} — 续} \\
\toprule
平台 & 编译 & Runtime \\
\midrule
\endhead
Linux x86\_64 Release & \texttt{-msha} 启用 intrinsic & CPUID + self-test \\
Linux x86\_64 Debug & 可能无 \texttt{-msha} & fallback Standard \\
Windows x64 MSVC & \texttt{\_M\_X64} + intrin & 同左 \\
ARM / RISC-V & 无 \texttt{EBBACKUP\_SHA\_NI\_COMPILED} & 恒 false \\
\bottomrule
\end{longtable}

CPUID 检测 EBX bit 29 对应 Intel SHA extensions；AMD Zen 部分型号亦支持。self-test 防止 errata 或虚拟化 CPU 误报。

\section{AES-GCM blob 布局与解密}

存储 blob 布局：\texttt{[nonce 12B][ciphertext N B][tag 16B]}。解密时 \texttt{Aes256GcmDecrypt} 校验 tag；失败返回 \texttt{Corrupt}。Windows BCrypt 路径失败时自动 fallback portable（与 encrypt 对称）。

\texttt{ChunkStore::DecodePayload} 在读取后根据 codec 加密位选择解密，再解压；verify 路径对解密明文重算 \texttt{ContentHash} 与 header hash 比对。

\section{ChunkRecordHeader CRC 计算}

v2 header CRC 算法：复制 header，\texttt{record\_crc32} 置零，对整个 48B 结构做 \texttt{Crc32}。v1 同理对 40B 计算。Open scan 时 CRC 不匹配尝试另一版本；均失败则视为 corrupt tail 起点。

\section{环境变量与运维调优汇总}

\begin{table}[htbp]
\centering
\caption{Digest 相关环境变量}
\begin{tabular}{@{}llp{6cm}@{}}
\toprule
变量 & 典型值 & 效果 \\
\midrule
\texttt{EBBACKUP\_DIGEST\_THREADS} & \texttt{4}（CI） & 固定 DigestPool worker 数 \\
 & 未设置 & \texttt{max(4, min(16, hw/2))} \\
 & \texttt{0} & 单线程（parity/debug） \\
\bottomrule
\end{tabular}
\end{table}

生产环境大核 CPU 可适当提高 \texttt{EBBACKUP\_DIGEST\_THREADS}（如 8--12），但超过 chunk 并行度收益递减；L5 profile 可观察 digest 阶段 \texttt{chunk\_ns} 占比。

\section{测试与 bench 覆盖}

\begin{table}[htbp]
\centering
\caption{Digest 相关测试矩阵（节选）}
\begin{tabular}{@{}llp{7cm}@{}}
\toprule
测试文件 & 覆盖点 \\
\midrule
\texttt{digest\_legacy\_test.cc} & Legacy SHA/HMAC 向量 \\
\texttt{digest\_standard\_test.cc} & Standard 路由、ContentHash \\
\texttt{digest\_sha\_ni\_test.cc} & SHA-NI parity、dispatch \\
\texttt{digest\_pool\_test.cc} & 多线程 HashRegions vs 串行 \\
\texttt{digest\_dual\_test.cc} & Legacy/Standard 仓库 round-trip \\
\texttt{pbkdf2\_test.cc} & PBKDF2 双栈一致 \\
\bottomrule
\end{tabular}
\end{table}

\section{本章小结}

Digest 子系统通过 \texttt{ContentHash} 统一路由 Legacy/Standard/SHA-NI，\texttt{DigestPool} 在 chunk 边界并行化摘要且可通过 \texttt{EBBACKUP\_DIGEST\_THREADS} 调优。加密在 store 之前应用，PBKDF2+salt 派生 CEK，但不改变 Content ID。下一章 \cref{ch:chunk-store} 描述 hash-keyed 持久化与 dedup 查重语义。
